\documentclass[10pt,a4paper]{amsart}
\usepackage[margin=19mm]{geometry}
\usepackage{amsmath,amssymb,mathtools}
\usepackage[scr=rsfs,cal=euler]{mathalfa}
\usepackage{xfrac}
\usepackage[unicode,hidelinks,bookmarks=false]{hyperref}
\newcommand{\ie}{i.e.\ }
\newcommand{\cf}{cf.\ }
\newcommand{\resp}{respectively\ }
\newcommand{\Subsection}{\subsection}
\providecommand{\nobreakdash}{\nobreak\hbox{-}}
\setlength{\emergencystretch}{2em}
\multlinegap0pt
\usepackage{stackengine}
\newtheorem{lemma}{Lemma}
\newcommand{\II}{{\mathcal I}}
\DeclareMathOperator{\adj}{adj}
\title{Homeostasis in Input-Output Networks: \texorpdfstring{\\} {} Structure, Classification and Applications}
\author{Fernando Antoneli, Martin Golubitsky, Jiaxin Jin, Ian Stewart}
\date{Source: arXiv:2405.03861v1 (CC BY 4.0)}
\begin{document}
\maketitle

\noindent\emph{Source and licence.} Homeostasis in Input-Output Networks: \texorpdfstring{\\} {} Structure, Classification and Applications. Version arXiv:2405.03861v1 (CC BY 4.0), distributed by arXiv under the Creative Commons Attribution 4.0 International licence, http://creativecommons.org/licenses/by/4.0/, as recorded in the arXiv licence metadata for this exact version and shown on the abstract page https://arxiv.org/abs/2405.03861v1. Full licence notice: \texttt{licenses/arxiv-2405.03861-CC-BY-4.0.txt}.\par\smallskip

\noindent\textit{Editorial scope.} Counted unit: the article's lemma t:z\_nabla (source0.tex lines 392--405). The article's complete original proof of it is delivered with it (source0.tex lines 407--428, one \texttt{proof} environment of 22 lines). No mathematics is written, repaired or re-derived: the statement and the proof are the article's own lines, in the article's own notation, and no retained line is substituted or re-flowed.  Retained uncounted as the source-local context the proof needs: lines 328--330.  Nothing was omitted from any retained span, so the package carries no omission record.  No retained line contains a citation, so the wrapper carries no bibliography and no bibliography entry.  No retained line contains a figure environment, an image-inclusion command or drawing code, so no image is delivered and the package declares no asset dependency.  Editorial formatting only, in addition: an AMS \texttt{amsart} preamble that restates the article's own declarations and macros used by the retained lines, namely the environments \texttt{lemma} on separate counters, together with the standard packages those lines need.  The retained environments print in this document as Lemma 1 in order of appearance, whereas the article prints the same items with the numbering of its own full document.  The complete native source of the article as distributed in the arXiv e-print for this exact version is bundled unchanged as \texttt{sources/158757/source0.tex}.
\begin{equation} \label{general_dynamics}
  \dot{X} = F(X, \II)
\end{equation}

\begin{lemma} \label{t:z_nabla}
Let $z(\mathcal{I})$ be an input-output function associated to a system of differential equations \eqref{general_dynamics}.
The gradient of $\nabla z(\II)$ with respect to the multiple input $\II$ is given by
\begin{equation} \label{e:z_prime}
\nabla z =  
\frac{-1}{\det J} 
\bigg(\langle \nabla \phi\, , \adj(J) F_{\II_M} \rangle \bigg)_{M=1}^N
\end{equation}
where $\adj(J)$ is the \emph{adjugate matrix} or \emph{classical adjoint} of $J$ (defined as the transpose of the cofactor matrix of $J$).
In particular, $\II_0$ is an infinitesimal homeostasis point of $z$ if and only if
\begin{equation} \label{e:cond_homeo}
 \langle \nabla \phi \, , \adj(J) F_{\II_M}\rangle \big|_{(X(\II_0),\II_0)} = 0 \;,\quad\forall\; M=1,\ldots,N
\end{equation}
\end{lemma}

\begin{proof}
The gradient of $z$ is
\begin{equation} \label{e:grad}
\nabla z = \bigg(\frac{\partial z}{\partial \II_1},\ldots,
\frac{\partial z}{\partial \II_N}\bigg)
\end{equation}
The partial derivative of $z(\mathcal{I})=\phi(X(\mathcal{I}))$ with respect to $\II_M$ (for each $M=1,\ldots,N$) is
\begin{equation} \label{e:chain_rule}
\frac{\partial z}{\partial \II_M} = \frac{\partial}{\partial \II_j} \phi(X(\II)) = \langle \nabla \phi\, ,  X_{\II_M} \rangle
\end{equation}
where $X_{\II_M}$ is the partial derivative of $X$ with respect to $\II_M$.
Implicit differentiation of the equation $F(X(\mathcal{I}),\mathcal{I})=0$, with respect to $\II_j$, yields the linear system
\begin{equation} \label{e:imp_diff}
J \, X_{\II_M} = - F_{\II_M}
\end{equation}
From \eqref{e:imp_diff} and the fact that $X(\II)$ is assumed to be a linearly stable equilibrium, for all $\II\in \Omega$, and hence $\det(J)\neq 0$ over $\Omega$, it follows that
\begin{equation} \label{e:adjoint}
X_{\II_M} = - \frac{1}{\det J} \adj(J) F_{\II_M}
\end{equation}
Therefore, substituting \eqref{e:adjoint} into \eqref{e:chain_rule} and then into \eqref{e:grad} we obtain \eqref{e:z_prime}.
\qed
\end{proof}

\end{document}
